%% ============================================================
%%  Rapport ML - Analyse Comportementale Clientèle Retail
%%  École Nationale d'Ingénieurs de Sfax — GI2 — 2025-2026
%% ============================================================
\documentclass[12pt, a4paper]{article}

%% ── Encodage & Langue ──────────────────────────────────────
\usepackage[utf8]{inputenc}
\usepackage[T1]{fontenc}
\usepackage[french]{babel}
\usepackage{lmodern}

%% ── Mise en page ───────────────────────────────────────────
\usepackage[top=2.5cm, bottom=2.5cm, left=2.5cm, right=2.5cm]{geometry}
\usepackage{microtype}
\usepackage{parskip}
\setlength{\parindent}{0pt}
\setlength{\parskip}{6pt}

%% ── Couleurs ───────────────────────────────────────────────
\usepackage{xcolor}
\definecolor{primary}{RGB}{44, 62, 80}      % bleu foncé
\definecolor{accent}{RGB}{52, 152, 219}     % bleu clair
\definecolor{success}{RGB}{39, 174, 96}     % vert
\definecolor{danger}{RGB}{231, 76, 60}      % rouge
\definecolor{lightgray}{RGB}{245, 245, 245}
\definecolor{codebg}{RGB}{248, 248, 248}

%% ── Graphiques & Figures ───────────────────────────────────
\usepackage{graphicx}
\usepackage{svg}
\usepackage{float}
\usepackage{caption}
\usepackage{subcaption}
\captionsetup{font=small, labelfont=bf, labelsep=period}

%% ── Tableaux ───────────────────────────────────────────────
\usepackage{booktabs}
\usepackage{tabularx}
\usepackage{array}
\usepackage{multirow}
\usepackage{colortbl}

%% ── Mise en forme des sections ─────────────────────────────
\usepackage{titlesec}
\titleformat{\section}
  {\large\bfseries\color{primary}}
  {\thesection}{1em}{}
  [\color{accent}\titlerule]
\titleformat{\subsection}
  {\normalsize\bfseries\color{primary}}
  {\thesubsection}{1em}{}
\titleformat{\subsubsection}
  {\normalsize\itshape\color{primary}}
  {\thesubsubsection}{1em}{}

%% ── En-têtes & Pieds de page ───────────────────────────────
\usepackage{fancyhdr}
\pagestyle{fancy}
\fancyhf{}
\fancyhead[L]{\small\color{primary}\textit{Analyse Comportementale Clientèle}}
\fancyhead[R]{\small\color{primary}\textit{Atelier Machine Learning}}
\fancyfoot[C]{\small\thepage}
\renewcommand{\headrulewidth}{0.4pt}
\renewcommand{\headrule}{\hbox to\headwidth{\color{accent}\leaders\hrule height \headrulewidth\hfill}}

%% ── Boîtes colorées ────────────────────────────────────────
\usepackage{tcolorbox}
\tcbuselibrary{skins, breakable}

\newtcolorbox{infobox}[1][]{
  colback=accent!8, colframe=accent, coltitle=white,
  fonttitle=\bfseries\small, title=#1,
  arc=4pt, boxrule=0.8pt, left=8pt, right=8pt,
  top=6pt, bottom=6pt
}

\newtcolorbox{codebox}{
  colback=codebg, colframe=primary!30,
  arc=3pt, boxrule=0.5pt, left=8pt, right=8pt,
  top=4pt, bottom=4pt, breakable
}

%% ── Code ───────────────────────────────────────────────────
\usepackage{listings}
\lstset{
  backgroundcolor=\color{codebg},
  basicstyle=\footnotesize\ttfamily,
  keywordstyle=\color{accent}\bfseries,
  commentstyle=\color{success}\itshape,
  stringstyle=\color{danger},
  breaklines=true,
  frame=single,
  framesep=4pt,
  rulecolor=\color{primary!30},
  numbers=left,
  numberstyle=\tiny\color{gray},
  numbersep=8pt,
  showstringspaces=false,
  tabsize=4,
  language=Python
}

%% ── Listes ─────────────────────────────────────────────────
\usepackage{enumitem}
\setlist[itemize]{leftmargin=1.5em, itemsep=2pt}
\setlist[enumerate]{leftmargin=1.5em, itemsep=2pt}

%% ── Liens ──────────────────────────────────────────────────
\usepackage{hyperref}
\hypersetup{
  colorlinks=true, linkcolor=primary,
  urlcolor=accent, citecolor=accent,
  pdftitle={Rapport ML - Analyse Comportementale Clientèle Retail},
  pdfauthor={Mahmoud Elloumi}
}

%% ── Commandes utiles ───────────────────────────────────────
\newcommand{\code}[1]{\texttt{\small#1}}
\newcommand{\feat}[1]{\textit{\textcolor{accent}{#1}}}

%% ============================================================
%%  DOCUMENT
%% ============================================================
\begin{document}

%% ── Page de garde ──────────────────────────────────────────
\thispagestyle{empty}
\begin{titlepage}

  % En-tête institutionnel
  \begin{minipage}[c]{0.30\textwidth}
    \centering\small
    République Tunisienne\\
    Ministère de l'Enseignement Supérieur,\\
    et de la Recherche Scientifique\\
    \textbf{Université de Sfax}\\
    École Nationale d'Ingénieurs de Sfax
  \end{minipage}
  \hfill
  \begin{minipage}[c]{0.30\textwidth}
    \centering
    \includesvg[width=0.55\textwidth]{Logo_ENIS.svg}
  \end{minipage}
  \hfill
  \begin{minipage}[c]{0.30\textwidth}
    \centering\small
    Cycle de Formation d'Ingénieurs\\
    dans la Discipline\\
    \textbf{Génie Informatique}
  \end{minipage}

  \vspace{4cm}
  \begin{center}
    {\Huge\bfseries\color{primary} COMPTE RENDU}\\[1.2cm]
    {\Huge\bfseries\color{primary} Atelier Machine Learning}\\[6pt]
    {\large\itshape\color{primary!70} Analyse Comportementale Clientèle Retail}\\[3cm]
    {\LARGE E-commerce de Cadeaux}\\[3cm]
    {\large\bfseries Préparé par \textsc{Mahmoud Elloumi}}\\[1.5cm]

    \rule{8cm}{0.6pt}\\[1cm]
    {\small\itshape Document pédagogique — Chaîne complète de traitement}\\[0.3cm]
    {\small Exploration $\rightarrow$ Préparation $\rightarrow$ Modélisation
            $\rightarrow$ Évaluation $\rightarrow$ \textbf{Déploiement}}\\[1.5cm]
    {\large\bfseries Année Universitaire : 2025-2026}
  \end{center}

\end{titlepage}

%% ── Table des matières ─────────────────────────────────────
\newpage
\tableofcontents
\newpage

%% ============================================================
\section{Introduction et Contexte}
%% ============================================================

Ce rapport s'inscrit dans le cadre du module de Machine Learning (GI2) de l'année
universitaire 2025-2026. Il détaille la conception et la mise en œuvre d'une chaîne
complète de traitement de données, de l'exploration initiale jusqu'au déploiement
d'une solution prédictive.

\subsection{Contexte de l'étude}

Nous intervenons en tant que Data Scientist pour une entreprise spécialisée dans le
\textbf{e-commerce de cadeaux}. Dans un secteur très concurrentiel, cette entreprise
souhaite tirer parti de ses données pour affiner sa connaissance de sa clientèle et
améliorer ses prises de décision.

\subsection{Objectifs métiers}

La mission principale est d'exploiter les données transactionnelles pour répondre
à trois enjeux stratégiques~:

\begin{itemize}
  \item \textbf{Personnalisation marketing :} Mieux comprendre les segments
        de clients afin de personnaliser efficacement les stratégies marketing.
  \item \textbf{Fidélisation (Anticipation du Churn) :} Identifier les signaux
        faibles permettant de réduire le taux de départ des clients (\textit{churn}).
  \item \textbf{Croissance :} Optimiser globalement le chiffre d'affaires.
\end{itemize}

\subsection{Présentation du jeu de données}

Nous nous appuyons sur la base \code{retail\_customers\_COMPLETE\_CATEGORICAL.csv},
qui regroupe \textbf{52 variables} issues de transactions réelles enrichies par des
données complémentaires. Cette base a été intentionnellement conçue pour être imparfaite
afin de simuler un vrai projet de Data Science.

%% ============================================================
\section{Configuration de l'environnement de travail}
%% ============================================================

\subsection{Arborescence et structure du projet}

Le projet respecte une architecture modulaire~:

\begin{codebox}
\begin{verbatim}
ML-Projet/
├── data/
│   ├── raw/          # Données brutes originales
│   ├── processed/    # Données nettoyées
│   └── train_test/   # Train / Test splittés
├── notebooks/        # Exploration Jupyter
├── src/
│   ├── preprocessing.py
│   ├── train_model.py
│   ├── predict.py
│   └── utils.py
├── models/           # Modèles sauvegardés (.pkl)
├── app/              # Application Flask
│   ├── app.py
│   ├── templates/
│   └── static/
├── reports/          # Visualisations & rapports
├── requirements.txt
└── README.md
\end{verbatim}
\end{codebox}

\subsection{Gestion de l'environnement virtuel}

Un environnement virtuel a été créé via~:

\begin{lstlisting}[language=bash]
python -m venv venv
venv\Scripts\Activate.ps1      # Windows PowerShell
pip install -r requirements.txt
\end{lstlisting}

Les dépendances sont figées dans \code{requirements.txt} (\code{pip freeze > requirements.txt})
pour assurer la reproductibilité.

%% ============================================================
\section{Analyse et Synthèse des Données}
%% ============================================================

\subsection{Exploration des features}

Le dataset contient \textbf{52 features} réparties en deux grandes familles~:

\begin{itemize}
  \item \textbf{Features numériques (1–34)~:} métriques RFM (\feat{Recency},
        \feat{Frequency}, \feat{MonetaryTotal}), comportement transactionnel
        (\feat{ReturnRatio}, \feat{CancelledTransactions}), profil client
        (\feat{Age}, \feat{SatisfactionScore}, \feat{SupportTicketsCount}).
  \item \textbf{Features catégorielles (35–52)~:} segments
        (\feat{RFMSegment}, \feat{SpendingCategory}, \feat{CustomerType}),
        temporelles et géographiques (\feat{FavoriteSeason}, \feat{Region},
        \feat{Country}), démographiques (\feat{AgeCategory}, \feat{Gender}).
\end{itemize}

\subsection{Identification des problèmes de qualité}

\begin{table}[H]
\centering
\caption{Problèmes de qualité identifiés dans le dataset}
\renewcommand{\arraystretch}{1.3}
\begin{tabularx}{\textwidth}{l l X l}
\toprule
\rowcolor{primary!90}
\color{white}\textbf{Type} &
\color{white}\textbf{Feature(s)} &
\color{white}\textbf{Description} &
\color{white}\textbf{Traitement} \\
\midrule
Valeurs manquantes  & \code{Age}                  & $\approx30\%$ de NaN               & Imputation médiane \\
\rowcolor{lightgray}
Valeurs aberrantes  & \code{SupportTicketsCount}  & Codes $-1$ et $999$                & Remplacement NaN \\
Valeurs aberrantes  & \code{SatisfactionScore}    & Codes $-1$, $0$ et $99$            & Remplacement NaN \\
\rowcolor{lightgray}
Formats inconsistants & \code{RegistrationDate}   & Formats UK et ISO mélangés         & Parsing dayfirst \\
Feature inutile     & \code{NewsletterSubscribed} & Constante ($100\%$ \textit{"Yes"}) & Suppression \\
\rowcolor{lightgray}
Données brutes      & \code{LastLoginIP}          & Identifiant unique                  & Suppression \\
Fuite de données    & \code{ChurnRiskCategory}    & Dérivé de Churn                    & Suppression \\
\rowcolor{lightgray}
Fuite de données    & \code{RFMSegment}           & Dérivé des features RFM            & Suppression \\
Déséquilibre classes & \code{Churn}               & $\approx33\%$ churners             & SMOTE \\
\bottomrule
\end{tabularx}
\end{table}

%% ============================================================
\section{Nettoyage et Ingénierie des Données}
%% ============================================================

\subsection{Traitement des valeurs aberrantes}

La fonction \code{handle\_aberrant\_values()} dans \code{utils.py} corrige~:

\begin{itemize}
  \item \feat{SupportTicketsCount}~: valeurs $-1$ et $999$ $\rightarrow$ \code{NaN}
        (130 occurrences corrigées)
  \item \feat{SatisfactionScore}~: valeurs $-1$, $0$ et $99$ $\rightarrow$ \code{NaN}
        (349 occurrences corrigées)
  \item \feat{MonetaryTotal}~: clip à $0$ pour éviter les montants négatifs extrêmes
\end{itemize}

\subsection{Feature Engineering}

De nouvelles variables ont été créées pour enrichir le pouvoir prédictif~:

\begin{lstlisting}[language=Python]
# Ratio monetaire / recence
df["MonetaryPerDay"] = df["MonetaryTotal"] / (df["Recency"] + 1)

# Panier moyen
df["AvgBasketValue"] = df["MonetaryTotal"] / (df["Frequency"] + 1)

# Ratio anciennete
df["TenureRatio"] = df["Recency"] / (df["CustomerTenureDays"] + 1)

# Frequence par jour
df["FrequencyPerDay"] = df["Frequency"] / (df["CustomerTenureDays"] + 1)

# Interaction RFM
df["Recency_x_Frequency"] = df["Recency"] * df["Frequency"]

# Log-transformations (distributions asymetriques)
df["LogMonetaryTotal"] = np.log1p(df["MonetaryTotal"].clip(lower=0))
df["LogFrequency"]     = np.log1p(df["Frequency"])
\end{lstlisting}

\subsection{Traitement des dates}

La variable \feat{RegistrationDate} présentait des formats hétérogènes
(UK~: \code{"12/03/10"}, ISO~: \code{"2010-03-12"}). Elle a été parsée avec
\code{dayfirst=True} puis décomposée en~: \feat{RegYear}, \feat{RegMonth},
\feat{RegDay}, \feat{RegDayOfWeek}, \feat{IsWeekendReg}.

%% ============================================================
\section{Transformation et Préparation à l'Apprentissage}
%% ============================================================

\subsection{Encodage des variables catégorielles}

Le pipeline de prétraitement distingue deux types de variables catégorielles~:

\begin{itemize}
  \item \textbf{Encodage Ordinal} (\code{OrdinalEncoder}) pour les variables
        à ordre naturel~:
        \begin{itemize}
          \item \feat{SpendingCategory}~: Low $<$ Medium $<$ High $<$ VIP
          \item \feat{LoyaltyLevel}~: Nouveau $<$ Jeune $<$ Établi $<$ Ancien
          \item \feat{AgeCategory}~: 18--24 $<$ 25--34 $<$ \ldots $<$ 65+
          \item \feat{BasketSizeCategory}~: Petit $<$ Moyen $<$ Grand
          \item \feat{PreferredTimeOfDay}~: Nuit $<$ Matin $<$ Midi $<$ Après-midi $<$ Soir
        \end{itemize}
  \item \textbf{One-Hot Encoding} (\code{OneHotEncoder}) pour les variables
        nominales sans ordre~: \feat{CustomerType}, \feat{FavoriteSeason},
        \feat{Region}, \feat{Gender}, \feat{AccountStatus}, \feat{Country}, etc.
\end{itemize}

\subsection{Séparation Train/Test et mise à l'échelle}

\begin{lstlisting}[language=Python]
X_train, X_test, y_train, y_test = train_test_split(
    X, y, test_size=0.2, random_state=42, stratify=y
)
# Fit UNIQUEMENT sur X_train -> eviter le Data Leakage
preprocessor.fit_transform(X_train)
preprocessor.transform(X_test)
\end{lstlisting}

Après prétraitement~: $X_{train}$ de dimension $(3\,497 \times 115)$ et
$X_{test}$ de dimension $(875 \times 115)$.

%% ============================================================
\section{Modélisation Avancée et Évaluation}
%% ============================================================

\subsection{Algorithmes et stratégie}

Quatre algorithmes ont été comparés pour ce problème de classification binaire~:

\begin{itemize}
  \item \textbf{Régression Logistique} — baseline linéaire interprétable
  \item \textbf{Random Forest} — ensemble basé sur le bagging
  \item \textbf{Gradient Boosting} — ensemble basé sur le boosting séquentiel
  \item \textbf{SVM (SVC)} — maximisation de la marge entre classes
\end{itemize}

Avant l'entraînement, la technique \textbf{SMOTE} a été appliquée sur le jeu
d'entraînement pour corriger le déséquilibre des classes~:

\begin{center}
\textit{Avant SMOTE}~: \{0~: 2\,334, 1~: 1\,163\}
\hspace{1cm}
\textit{Après SMOTE}~: \{0~: 2\,334, 1~: 2\,334\}
\end{center}

\subsection{Optimisation des hyperparamètres (GridSearchCV)}

Une recherche exhaustive a été effectuée pour chaque algorithme via
\code{GridSearchCV} avec \textbf{cv=5} et scoring \code{roc\_auc}~:

\begin{itemize}
  \item Random Forest~: \code{n\_estimators} $\in \{100, 200\}$,
        \code{max\_depth} $\in \{\text{None}, 10, 20\}$
  \item Gradient Boosting~: \code{n\_estimators} $\in \{100, 200\}$,
        \code{learning\_rate} $\in \{0.05, 0.1\}$
  \item Régression Logistique~: \code{C} $\in \{0.01, 0.1, 1, 10\}$
\end{itemize}

\subsection{Résultats comparatifs}

\begin{table}[H]
\centering
\caption{Comparaison des performances des modèles sur le jeu de test (20\,\%)}
\renewcommand{\arraystretch}{1.4}
\begin{tabularx}{\textwidth}{l c c c c c}
\toprule
\rowcolor{primary!90}
\color{white}\textbf{Modèle} &
\color{white}\textbf{Accuracy} &
\color{white}\textbf{Precision} &
\color{white}\textbf{Recall} &
\color{white}\textbf{F1-Score} &
\color{white}\textbf{AUC-ROC} \\
\midrule
Régression Logistique & 0.9943 & 0.99 & 0.99 & 0.99 & 0.9984 \\
\rowcolor{lightgray}
\textbf{Random Forest} & \textbf{1.0000} & \textbf{1.00} & \textbf{1.00} & \textbf{1.00} & \textbf{1.0000} \\
Gradient Boosting     & 1.0000 & 1.00 & 1.00 & 1.00 & 1.0000 \\
\bottomrule
\end{tabularx}
\end{table}

Le \textbf{Random Forest} (200 arbres, profondeur illimitée) a été sélectionné
comme modèle final.

\subsection{Visualisations des performances}

\begin{figure}[H]
  \centering
  \begin{subfigure}[b]{0.48\textwidth}
    \centering
    \includegraphics[width=\textwidth]{confusion_matrix_Random_Forest.png}
    \caption{Matrice de confusion — Random Forest}
  \end{subfigure}
  \hfill
  \begin{subfigure}[b]{0.48\textwidth}
    \centering
    \includegraphics[width=\textwidth]{feature_importance_Random_Forest.png}
    \caption{Top 20 features les plus importantes}
  \end{subfigure}
  \caption{Évaluation du modèle Random Forest sur le jeu de test}
\end{figure}

%% ============================================================
\section{Déploiement~: Interface avec Flask}
%% ============================================================

\subsection{Architecture de l'application}

L'application web a été développée avec le micro-framework \textbf{Flask}.
Elle expose deux routes principales~:

\begin{itemize}
  \item \code{GET /}~: page d'accueil avec upload CSV
  \item \code{POST /predict}~: traitement et affichage des prédictions
\end{itemize}

Le pipeline d'inférence charge les artefacts sauvegardés
(\code{preprocessor.pkl} et \code{churn\_model.pkl}) et applique exactement
la même transformation que lors de l'entraînement.

\subsection{Page d'accueil — Upload CSV}

La page d'accueil permet d'importer un fichier CSV contenant les données
clients. Toutes les colonnes du dataset original sont supportées.

\begin{figure}[H]
  \centering
  
\includegraphics[width=0.85\textwidth]{screenshot_index.png}
  \caption{Interface d'accueil — Import de fichier CSV}
\end{figure}

\subsection{Tableau de bord analytique — Statistiques en ligne}

Après traitement, la page de résultats affiche~:

\begin{itemize}
  \item \textbf{KPI Cards}~: nombre de clients analysés, churners détectés,
        taux de churn prédit, risque moyen
  \item \textbf{Graphique Donut}~: répartition Fidèles / Churners
  \item \textbf{Graphique Pie}~: distribution des 4 niveaux de risque
        (Critique / Élevé / Modéré / Faible)
  \item \textbf{Histogramme}~: nombre de clients par tranche de probabilité
        (0--10\,\%, 10--20\,\%, \ldots)
  \item \textbf{Barres horizontales}~: Top 10 clients les plus à risque
\end{itemize}

\begin{figure}[H]
  \centering
  \includegraphics[width=\textwidth]{screenshot_stats.png}
  \caption{Tableau de bord — KPI et graphiques interactifs (Chart.js)}
\end{figure}

\begin{figure}[H]
  \centering
  \begin{subfigure}[b]{0.48\textwidth}
    \centering
    \includegraphics[width=\textwidth]{screenshot_charts1.png}
    \caption{Donut et distribution des risques}
  \end{subfigure}
  \hfill
  \begin{subfigure}[b]{0.48\textwidth}
    \centering
    \includegraphics[width=\textwidth]{screenshot_charts2.png}
    \caption{Histogramme et Top 10 clients}
  \end{subfigure}
  \caption{Graphiques statistiques générés automatiquement}
\end{figure}

\subsection{Tableau détaillé des prédictions}

Chaque client est affiché avec sa prédiction, sa probabilité de départ
(barre de progression colorée) et son niveau de risque interprété.

\begin{figure}[H]
  \centering
  \includegraphics[width=\textwidth]{screenshot_table.png}
  \caption{Tableau des prédictions — Trié par risque décroissant}
\end{figure}

%% ============================================================
\section{Conclusion et Recommandations}
%% ============================================================

\subsection{Bilan technique}

Ce projet a couvert l'intégralité du cycle d'un projet de Data Science~:

\begin{infobox}[Résumé de la chaîne de traitement]
\begin{itemize}
  \item \textbf{Données}~: 4\,372 clients $\times$ 52 features brutes
        $\rightarrow$ 115 features après preprocessing
  \item \textbf{Nettoyage}~: correction de 479 valeurs aberrantes,
        5 colonnes supprimées (fuite/inutiles)
  \item \textbf{Encodage}~: OrdinalEncoder (5 variables) +
        OneHotEncoder (nominales)
  \item \textbf{Équilibrage}~: SMOTE (1\,163 $\rightarrow$ 2\,334 churners)
  \item \textbf{Modèle final}~: Random Forest — AUC-ROC = 1.000
  \item \textbf{Déploiement}~: API Flask + dashboard interactif Chart.js
\end{itemize}
\end{infobox}

\subsection{Recommandations métier}

Les résultats obtenus permettent de proposer les actions suivantes~:

\begin{itemize}
  \item \textbf{Anticipation du churn}~: identifier en temps réel les clients
        à risque (probabilité $\geq 75\,\%$) et déclencher des actions de
        rétention ciblées (réductions, offres personnalisées).
  \item \textbf{Personnalisation marketing}~: adapter les campagnes selon les
        features les plus influentes (ancienneté, fréquence, montant moyen).
  \item \textbf{Optimisation des ressources}~: concentrer les efforts de
        fidélisation sur les clients à forte valeur monétaire identifiés
        comme à risque élevé.
\end{itemize}

\end{document}

